%% TRANSCRIBED (PDF 154-162 = printed IV-9..IV-17)
%% Sec. 2 "The Galois modules attached to E".
%% NOTE: the tail of Sec. 1.4 (the last two paragraphs before the Sec. 2
%% heading, printed on IV-9 = PDF 154) is transcribed at the end of
%% ch4_s1.tex; this file starts at the Sec. 2 heading.
%% NOTE (p. IV-17): the printed text reads "... are open in Aut(V) by the
%% theorem in section 4".  Chapter IV has only §§1-3; the reference is to
%% the theorem of 2.2 (g_l = End(V_l)), i.e. a misprint for "section 2".
%% Kept as printed (see STATE.md).

\sechead{2}{The Galois modules attached to $E$}{THE GALOIS MODULES ATTACHED TO $E$}

In this section, $E$ denotes an elliptic curve over $K$.  We are
interested in the structure of the Galois modules $E_{\ell^n}$, $T_\ell$,
$V_\ell$ defined in \S\ref{sec:IV.1.3}.

\subsechead{2.1}{The irreducibility theorem}

Recall first that the ring $\End_K(E)$ of $K$-endomorphisms of $E$ is
either $\Z$ or of rank $2$ over $\Z$.  In the first case, we say that $E$
has ``no complex multiplication over $K$.''  If the same is true for any
finite extension of $K$, we say that $E$ has ``no complex
multiplication.''

\begin{theorem}\label{thm:IV.2.1.1}
Assume that $E$ has no complex multiplication over $K$.  Then:
\begin{alist}
\item $V_\ell$ is irreducible for all primes $\ell$;
\item $E_\ell$ is irreducible for almost all primes $\ell$.
\end{alist}
\end{theorem}

We need the following elementary result:

\begin{lemma}\label{lem:IV.2.1.1}
Let $E$ be an elliptic curve defined over $K$ with $\End_K(E) = \Z$.
Then, if $E' \to E$, $E'' \to E$ are $K$-isogenies with non-isomorphic
cyclic kernels, the curves $E'$ and $E''$ are non-isomorphic over $K$.
\end{lemma}

\begin{proof}
Let $n'$ and $n''$ be respectively the orders of the kernels of
$E' \to E$ and $E'' \to E$.  Suppose that $E'$ and $E''$ are isomorphic
over $K$, and let $E' \to E''$ be an isomorphism.  If $E \to E'$ is the
transpose of the isogeny $E' \to E$, it has a cyclic kernel of order
$n'$, and hence the isogeny $E \to E$, obtained by composition of
$E \to E'$, $E' \to E''$, $E'' \to E$, has for kernel an extension of
$\Z/n''\Z$ by $\Z/n'\Z$.  But, since $\End_K(E) = \Z$, this isogeny must
be multiplication by an integer $a$, and its kernel must therefore be of
the form $\Z/a\Z \times \Z/a\Z$.  Hence $n'$ and $n''$ divide $a$.
Since $a^2 = n'n''$, we obtain $a = n' = n''$, a contradiction.
\end{proof}

\runhead{Proof of the theorem}

(a) It suffices to show that, if $\End_K(E) = \Z$, there is no
one-dimensional $\Q_\ell$-subspace of $V_\ell$ stable under
$\Gal(\bar{K}/K)$.  Suppose there were one; its intersection $X$ with
$T_\ell$ would be a submodule of $T_\ell$ with $X$ and $T_\ell/X$ free
$\Z_\ell$-modules of rank $1$.  For $n \geq 0$, consider the image $X(n)$
of $X$ in $E_{\ell^n} = T/\ell^n T$.  This is a submodule of
$E_{\ell^n}$ which is cyclic of order $\ell^n$ and stable by
$\Gal(\bar{K}/K)$.  Hence it corresponds to a finite $K$-algebraic
subgroup of $E$ and one can define the quotient curve
$E(n) = E/X(n)$.  The kernel of the isogeny $E \to E(n)$ is cyclic of
order $\ell^n$.  The above lemma then shows that the curves $E(n)$,
$n \geq 0$, are pairwise non-isomorphic, contradicting the corollary to
\v{S}afarevi\v{c}'s theorem (\S\ref{sec:IV.1.4}).

(b) If $E_\ell$ is not irreducible, there exists a Galois submodule
$X_\ell$ of $E$ which is one-dimensional over $\F_\ell$.  In the same way
as above, this defines an isogeny $E \to E/X_\ell$ whose kernel is cyclic
of order $\ell$.  The above lemma shows that the curves which correspond
to different values of $\ell$ are non-isomorphic, and one again applies
the corollary to \v{S}afarevi\v{c}'s theorem.

\runhead{Remark}

One can prove part (a) of the above theorem by a quite different method
(cf.~\cite{serre-lieg}, \S3.4); instead of the \v{S}afarevi\v{c}'s
theorem, one uses the properties of the decomposition and inertia
subgroups of $\Ima(\rho_\ell)$, cf.~Appendix.

\subsechead{2.2}{Determination of the Lie algebra of $G_\ell$}

Let $G_\ell = \Ima(\rho_\ell)$ denote the image of $\Gal(\bar{K}/K)$ in
$\Aut(T_\ell)$, and let $\mathfrak{g}_\ell \subset \End(V_\ell)$ be the
Lie algebra of $G_\ell$.

\begin{theorem}\label{thm:IV.2.2.1}
If $E$ has no complex multiplication (cf.~\ref{sec:IV.2.1}), then
$\mathfrak{g}_\ell = \End(V_\ell)$, i.e.~$G_\ell$ is open in
$\Aut(T_\ell)$.
\end{theorem}

\begin{proof}
The irreducibility theorem of \ref{sec:IV.2.1} shows that, for any open subgroup $U$
of $G_\ell$, $V_\ell$ is an irreducible $U$-module.  Hence, $V_\ell$ is
an irreducible $\mathfrak{g}_\ell$-module.  By Schur's lemma, it follows
that the commuting algebra $\mathfrak{g}'_\ell$ of $\mathfrak{g}_\ell$ in
$\End(V_\ell)$ is a field; since $\dim V_\ell = 2$, this field is either
$\Q_\ell$ or a quadratic extension of $\Q_\ell$.  If
$\mathfrak{g}'_\ell = \Q_\ell$, then $\mathfrak{g}_\ell$ is equal to
either $\End(V_\ell)$, or the subalgebra $\mathfrak{sl}(V_\ell)$ of
$\End(V_\ell)$ consisting of the endomorphisms with trace $0$; but, in
the second case, the action of $\mathfrak{g}_\ell$ on $\Lambda^2 V_\ell$
would be trivial, and this would contradict the fact that the Galois
modules $\Lambda^2 V_\ell$ and $V_\ell(\mu)$ are isomorphic
(chap.~\ref{chap:I}, \S\ref{sec:I.1.2}).  Hence
$\mathfrak{g}_\ell = \mathfrak{sl}(V_\ell)$ is impossible.

Suppose now that $\mathfrak{g}'_\ell$ is a quadratic extension of
$\Q_\ell$.  Then $V_\ell$ is a one-dimensional
$\mathfrak{g}'_\ell$-vector space and the commuting algebra of
$\mathfrak{g}'_\ell$ in $\End(V_\ell)$ is $\mathfrak{g}'_\ell$ itself.
Hence $\mathfrak{g}_\ell$ is contained in $\mathfrak{g}'_\ell$, and is
abelian ($\mathfrak{g}'_\ell$ is a ``non-split Cartan algebra'' of
$\End(V_\ell)$).  After replacing $K$ by a finite extension (this does
not affect $\mathfrak{g}_\ell$, cf.~chap.~\ref{chap:I},
\S\ref{sec:I.1.1}), we may then suppose that $G_\ell$ itself is abelian.
The \textit{l}-adic representation $V_\ell$ is then semi-simple, abelian
and rational.  It is, moreover, locally algebraic.  To see this, we first
remark that, at a place $v$ dividing $\ell$, we have $v(j) \geq 0$ since
otherwise the decomposition group of $v$ in $G_\ell$ would be non-abelian
by Tate's theory (cf.~Appendix, \S\ref{sec:IV.A.1.3}); hence, after a
finite extension of $K$, we can assume that $E$ has good reduction at all
places $v$ dividing $\ell$ (cf.~\ref{sec:IV.1.2}).  Let $E(\ell)$ be the
$\ell$-divisible group attached to $E$ at $v$ (cf.~Tate
\cite{tate-pdiv}, 2.1, example (a)).  We have
$V_\ell \sim V_\ell(E(\ell))$ and this module is known to be of
Hodge-Tate type (\emph{loc.~cit.}, \S4).  Using another result of Tate
(chap.~\ref{chap:III}, \S\ref{sec:III.1.2}), this implies that the
representation $V_\ell$ is locally algebraic, as claimed above.  (This
could also be seen by using, instead of the theory of Hodge-Tate modules,
the local results of the Appendix, \S\ref{sec:IV.A.2}.)

We may now apply to $V_\ell$ the results of chap.~\ref{chap:III},
\S\ref{sec:III.2.3}.  Hence, there is, for each prime $\ell'$, a
rational, abelian, semi-simple $\ell'$-adic representation $W_{\ell'}$
compatible with $V_\ell$.  But $V_{\ell'}$ is compatible with $V_\ell$,
and $V_{\ell'}$ is semi-simple.  Hence $V_{\ell'}$ is isomorphic to
$W_{\ell'}$ (cf.~chap.~\ref{chap:I}, \S\ref{sec:I.2.3}).  But we know
(chap.~\ref{chap:III}, \S\ref{sec:III.2.3}) that we may choose $\ell'$
such that $W_{\ell'}$ is the direct sum of one-dimensional subspaces
stable under $\Gal(\bar{K}/K)$.  This contradicts the irreducibility of
$V_{\ell'}$.  Hence, we must have $\mathfrak{g}'_\ell = \Q_\ell$ and
$\mathfrak{g}_\ell = \End(V_\ell)$, q.e.d.
\end{proof}

\runhead{Remark}

If $E$ has complex multiplication, and
$L = \Q \tens \End(E \times_K \bar{K})$ is the corresponding imaginary
quadratic field, one shows easily that $\mathfrak{g}_\ell$ is the Cartan
subalgebra of $\End(V_\ell)$ defined by $L_\ell = \Q_\ell \tens L$.  It
splits if and only if $\ell$ decomposes in $L$.

\runhead{Exercises}

(In these exercises, we assume $E$ has no complex multiplication.  Let
$S$ be the set of places $v \in \Sigma_K$ where $E$ has bad reduction.
If $v \in \Sigma_K - S$, we denote by $F_v$ the Frobenius endomorphism of
the reduced curve $\tilde{E}_v$; if $\ell \neq p_v$, we identify $F_v$ to
the corresponding automorphism of $T_\ell$.)

\begin{nlist}
\item Let $H(X,Y)$ be a polynomial in two indeterminates $X$, $Y$ with
coefficients in a field of characteristic zero.  Let $V_H$ be the set of
those $v \in \Sigma_K - S$ for which
$H(\mathrm{Tr}(F_v), \mathrm{N}v) = 0$.  If $H$ is not the zero
polynomial, show that $V_H$ has density $0$.  (Show that the set of
$g \in \GL(2,\Z_\ell)$ with
$H(\mathrm{Tr}(g), \det(g)) = 0$ has Haar measure zero.)

\item The eigenvalues of $F_v$ may be identified with complex numbers of
the form
\[
\sqrt{\mathrm{N}v}\, e^{\pm i\varphi_v}, \qquad 0 < \varphi_v < \pi
\]
cf.~chap.~\ref{chap:I}, Appendix \S\ref{sec:I.A.2}.  Show that the set of
$v$ for which $\varphi_v$ is a given angle $\varphi$ has density zero.
(Show that $\mathrm{Tr}(F_v)^2 = 4(\mathrm{N}v)\cos^2 \varphi_v$ and then
use the preceding exercise.)

\item Let $L_v = \Q(F_v)$ be the field generated by $F_v$.  By the
preceding exercise, $L_v$ is quadratic imaginary except for a set of $v$
of density $0$.
\begin{alist}
\item Let $\ell$ be a fixed prime.  Let $C$ be a semi-simple commutative
$\Q_\ell$-algebra of rank $2$.  Let $X_C$ be the set of elements
$s \in \Aut(V_\ell)$ such that the subalgebra $\Q_\ell[s]$ of
$\End(V_\ell)$ generated by $s$ is isomorphic to $C$.  Show that $X_C$ is
open in $\Aut(V_\ell)$, and show that it has a non-empty intersection
with every open subgroup of $\Aut(V_\ell)$, in particular, with $G_\ell$.

\item Show that $F_v \in X_C$ if and only if the field $L_v$ is quadratic
and $L_v \tens \Q_\ell$ is isomorphic to $C$.

\item Let $\ell_1, \ldots, \ell_n$ be distinct prime numbers, and choose
for each an algebra $C_i$ of the type considered in (a).  Show that the
set of $v$ for which $F_v \in X_{C_i}$ for $i = 1, \ldots, n$ has density
$> 0$.  (Use the fact that the image of $\Gal(\bar{K}/K)$ in any finite
product of the $\Aut(V_\ell)$ is open; this is an easy consequence of the
theorem proved above.)

\item Deduce that, for any finite set $P$ of prime numbers, there exist an
infinity of $v$ such that $L_v$ is ramified at all $\ell \in P$.  In
particular, there are an infinite number of distinct fields $L_v$.
\end{alist}
\end{nlist}

\subsechead{2.3}{The isogeny theorem}

\begin{theorem}\label{thm:IV.2.3.1}
Let $E$ and $E'$ be elliptic curves over $K$, let $\ell$ be a prime
number and let $V_\ell(E)$ and $V_\ell(E')$ be the corresponding
$\ell$-adic representations of $K$.  Suppose that the Galois modules
$V_\ell(E)$ and $V_\ell(E')$ are isomorphic and that the modular
invariant $j$ of $E$ (cf.~\ref{sec:IV.1.1}) is not an integer of $K$.  Then $E$ and
$E'$ are $K$-isogenous.
\end{theorem}

We need the following result:

\begin{proposition}\label{prop:IV.2.3.1}
Let $E$ and $E'$ be elliptic curves over $K$.  The following conditions
are equivalent:
\begin{alist}
\item The Galois modules $V_\ell(E)$ and $V_\ell(E')$ are isomorphic for
all $\ell$.
\item The Galois modules $V_\ell(E)$ and $V_\ell(E')$ are isomorphic for
one $\ell$.
\item If $F_v$ and $F'_v$ are the Frobeniuses of the reduced curves
$\tilde{E}_v$ and $\tilde{E}'_v$, we have
$\mathrm{Tr}(F_v) = \mathrm{Tr}(F'_v)$ for all $v$ where there is good
reduction.
\item For a set of places of $K$ of density one we have
$\mathrm{Tr}(F_v) = \mathrm{Tr}(F'_v)$.
\end{alist}
\end{proposition}

Clearly (a) implies (b), and (c) implies (d).  The implication
(b) $\Rightarrow$ (c) follows from the fact that $\mathrm{Tr}(F_v)$ is
known when $V_\ell$ is known.  To prove (d) $\Rightarrow$ (a) one remarks
first that the representations of $\Gal(\bar{K}/K)$ in $V_\ell(E)$ and
$V_\ell(E')$ have the same trace, by \v{C}ebotarev's density theorem
(chap.~\ref{chap:I}, \S\ref{sec:I.2.2}).  Moreover, $V_\ell(E)$ (and also
$V_\ell(E')$) is \emph{semi-simple}.  This is clear if $E$ has no complex
multiplication over $K$ since $V_\ell(E)$ is then irreducible (\ref{sec:IV.2.1}); if
$E$ has complex multiplication, it follows from the Remark in \ref{sec:IV.2.2}.  Since
$V_\ell(E)$ and $V_\ell(E')$ are semi-simple and have the same trace,
they are isomorphic.

\runhead{Remarks}

\begin{nlist}
\item If $E$ and $E'$ have good reduction at $v$, let $t_v$
(resp.~$t'_v$) be the number of $k_v$-points of $\tilde{E}_v$
(resp.~$\tilde{E}'_v$).  We have the formulas (cf.~\ref{sec:IV.1.3}):
\[
t_v = 1 - \mathrm{Tr}(F_v) + \mathrm{N}v
\]
\[
t'_v = 1 - \mathrm{Tr}(F'_v) + \mathrm{N}v
\]
Hence condition (c) (resp.~condition (d)) is equivalent to saying that
$t_v = t'_v$ for all $v$ where there is good reduction (resp.~for a set
of $v$'s of density one).

\item If $E$ and $E'$ are $K$-isogenous, it is clear that conditions
(a), (b), (c), (d) are satisfied.
\end{nlist}

\begin{proof}[Proof of the theorem]
In view of Remark 2) above, it suffices to show that the equivalent
conditions (a), (b), (c), (d) imply that the elliptic curves $E$ and $E'$
are isogenous when the modular invariant $j$ of $E$ is not an integer of
$K$.  Let $v$ be a place of $K$ such that $v(j) < 0$, and let $p$ be the
characteristic of the residue field $k_v$.

If $j' = j(E')$, we first show that $v(j')$ is also $< 0$.  Suppose that
$v(j') \geq 0$.  Then, after possibly replacing $K$ by a finite extension,
we may assume that $E'$ has good reduction at $v$.  Then, if
$\ell \neq p$, the Galois-module $V_\ell(E')$ is unramified at $v$
(cf.~\ref{sec:IV.1.3}); but $V_\ell(E)$ is ramified at $v$: this follows either from
the criterion of N\'eron--Ogg--\v{S}afarevi\v{c} (\ref{sec:IV.1.3}) or from the
determination of the inertia group given in the Appendix,
\S\ref{sec:IV.A.1.3}.  This contradicts the fact that $V_\ell(E)$ and
$V_\ell(E')$ are isomorphic.

Let now $q$ and $q'$ be the elements of $K_v$ which correspond to $j$ and
$j'$ in Tate's theory (cf.~Appendix \S\ref{sec:IV.A.1.1}), and let $E_q$
and $E_{q'}$ be the corresponding elliptic curves (\emph{loc.~cit.}).
Since $E$ and $E_q$ have the same modular invariant $j$, there is a
finite extension $K'$ of $K_v$ where they become isomorphic, and we can
do the same for $E'$ and $E_{q'}$.  Hence, the Tate modules $T_p(E_q)$
and $T_p(E_{q'})$ become isomorphic over $K'$.  But, in this case the
isogeny theorem is true (cf.~Appendix \S\ref{sec:IV.A.1.4}), i.e.~the
curves $E_q$ and $E_{q'}$, hence also $E$ and $E'$, are $K'$-isogenous.
However, if two elliptic curves are isogenous over some extension of the
ground field, they are isogenous over a \emph{finite} extension of the
ground field.  We may thus choose a finite extension $L$ of $K$ and an
$L$-isogeny $f\colon E \times_K L \to E' \times_K L$.  We will show that
$f$ is automatically defined over $K$.  For this, it suffices to show
that $f = {}^s f$ for all $s \in \Gal(\bar{K}/K)$, or, equivalently, that
$V(f)\colon V_p(E) \to V_p(E')$ commutes with the action of Galois.
However, if $G_L = \Gal(\bar{K}/L)$ is the open subgroup of
$G = \Gal(\bar{K}/K)$ which corresponds to $L$, then $V(f)$ commutes with
the action of $G_L$.  It is then enough to show that
$\mathrm{Hom}_{G_L}(V,V') = \mathrm{Hom}_G(V,V')$.  But $V$ and $V'$ are
isomorphic as $G$-modules.  Hence we have to show that
$\mathrm{End}_{G_L}(V) = \mathrm{End}_G(V)$.  But this is clearly true; in
fact, $G$ and $G_L$ are open in $\Aut(V)$ by the theorem in section 4,
and hence their commuting algebra is reduced to the homotheties in each
case, i.e.~$\mathrm{End}_{G_L}(V) = \mathrm{End}_G(V) = \Q_p$.  This
completes the proof of the theorem.
\end{proof}

\runhead{Remark}

It is very likely that the theorem is true without the hypothesis that
$j$ is not integral.  This could be proved (by Tate's method
\cite{tate-endo}) if the following generalization of \v{S}afarevi\v{c}'s
theorem were true: given a finite subset $S$ of $\Sigma_K$, the abelian
varieties over $K$, of dimension $2$, with polarization of degree one,
and good reduction outside $S$, are in finite number (up to isomorphism).
$\longrightarrow$ [this has been proved by Faltings, see
\cite{faltings-endl}, \cite{faltings-wustholz}, \cite{szpiro}.]
